\let\textcircled=\pgftextcircled
\chapter{Pr\'esentation du Laboratoire de Sciences du Numérique de Nantes}
\label{chap:structure}

\textit{\initial{D}ans ce chapitre, on va présenter la structure où on a eu à effectuer notre stage, le Laboratoire des Sciences du Numérique de Nantes en abrégé LS2N : }

\mtcselectlanguage{french}
\minitoc

\newpage    

\section{Présentation du LS2N}
Le Laboratoire des Sciences du Numérique de Nantes (LS2N) a été créé en janvier 2017 afin de répondre à l’ambition de faire progresser significativement la visibilité de la recherche en sciences du numérique à Nantes. Le LS2N résulte essentiellement de la fusion des ex-UMR IRCCyN (Institut de Recherche en Communications et Cybernétique de Nantes) et LINA (Laboratoire d’Informatique de Nantes Atlantique). Cette jeune unité mixte de recherche met en commun les forces de recherche dans le domaine du numérique de trois établissements d’enseignement supérieur (Université de Nantes, Centrale Nantes, l’IMT Atlantique/campus de Nantes), du CNRS et de l’Inria. Le LS2N est la plus grosse unité de recherche publique sur le site de Nantes et en région Pays de la Loire.

L’effectif du LS2N s’élève à plus de 500 personnes dont la moitié de permanents, répartis sur 5 sites : Faculté des Sciences et Techniques, Centrale Nantes, Polytech Nantes, IMT Atlantique et IUT de Nantes. La recherche du LS2N est effectuée au sein de 22 équipes de recherche, structurées autour de 5 grands pôles scientifiques : Signaux, Images, Ergonomie et Langues ; Sciences des Données et de la Décision ; Science du Logiciel et des Systèmes Distribués ; Conception et Conduite de Systèmes et Robotique, Procédés, Calcul. À ces 5 pôles s’ajoutent les 6 thèmes transverses d’application : Industrie et entreprise du futur, Énergie et impacts environnementaux, Sciences du vivant, Véhicules et mobilités, Cultures numériques et Technologie Numérique pour l’Éducation Ouverte.

\section{Organisation}
Le laboratoire est présent sur 5 sites de la métropole nantaise :
\begin{itemize}
    \item Le campus de Centrale Nantes (CN) (176 personnes),
    \item Le campus de la faculté des sciences (FST) (103 personnes),
    \item Le campus de l’IMT Atlantique (IMT-A) (90 personnes),
    \item Le campus de l’école polytechnique (Polytech) (55 personnes),
    \item Le campus de l’Institut Universitaire de Technologie (IUT) (26 personnes).
\end{itemize}

\begin{figure}[H]
    \centering
    \includegraphics[width=0.5\linewidth]{fig02/planSiteNantes.png}
    \caption{Plan des sites du LS2N à Nantes}
\end{figure}

